\vfill%
\pagebreak

\section*{Exercises}
\markright{\MakeUppercase{Exercises}}%
\subsection*{Exercise 1}

Write a function \token{smallest} that returns the smaller of two \token{i32}.
Then overload it with a version that takes three, written with calls to the
first, and no \token{if} of its own.

\subsection*{Exercise 2}

What does the following program print? Work it out by hand, then check your
answer by compiling and running it.

\begin{lstlisting}[style=coloredverbatim]
use std::io;

fn show(label: [c8], value: i32, unit: [c8] = "") {
  println(label, ": ", value, unit);
}

fn twice(x: i32)-> i32 {
  x * 2
}

fn main() {
  let n = 5;
  show("n", n);
  show("twice", n.twice(), unit-> " points");
  show(value-> n.twice().twice(), "four times");
  show(unit-> "!", "sum", twice(n) + n);
}
\end{lstlisting}

\subsection*{Exercise 3}

The following program should print the average of 4 and 10, kept within the
bounds 0 and 5. It contains three errors. Find them, correct them, and compile the
program to test your corrections.

\begin{lstlisting}[style=coloredverbatimError, caption=Ymir program with errors]
use std::io;

fn average(a: i32, b: i32)-> i32 {
  (a + b) / 2;
}

fn clamp(value: i32, @low: i32, @high: i32)-> i32 {
  if value < low {
    value = low;
  }
  if value > high { high } else { value }
}

fn main() {
  let m = average(4, 10);
  println(clamp(m, 0, 5));
}
\end{lstlisting}

\subsection*{Exercise 4}

Write a recursive function \token{power(base: i64, exponent: i32)-> i64} that
computes \token{base} to the power \token{exponent}, knowing that any number to
the power 0 is 1, and that $b^e = b \times b^{e - 1}$. What happens if it is
called with a negative exponent? Change the function so that it returns 1
instead.

\subsection*{Exercise 5}

Exercise~4 of Chapter~\ref{chap:control} computed the sum of the digits of a
number with a loop. Write it as a recursive function instead. (\textit{Hint: a
number below 10 is its own sum of digits; for a larger one, what is the relation
between the sum of the digits of \token{n} and that of \token{n / 10}?})

\subsection*{Exercise 6}

The Fibonacci sequence starts with 0 and 1, and each number after them is the sum
of the two before it: 0, 1, 1, 2, 3, 5, 8, 13, \dots\ Write a recursive function
\token{fibonacci(n: i64)-> i64} that returns the number at position \token{n},
counting from 0, straight from that definition. Run it for 30, 40 and 45. Why
does it get so slow, when a loop adds up millions of numbers in the blink of an
eye? Write a version with a loop that answers at once, even for 90.

\subsection*{Exercise 7}

The greatest common divisor of two numbers is the largest number that divides
both. Euclid found, some 2\,300 years ago, that the greatest common divisor of
\token{a} and 0 is \token{a}, and that the greatest common divisor of \token{a}
and \token{b} is that of \token{b} and \token{a \% b}. Write a recursive function
\token{gcd} from these two rules, and check that the greatest common divisor of
1071 and 462 is 21. Why does the recursion always end?

\subsection*{Exercise 8}

Following the rules of Section~\ref{sec:functions:ackermann}, compute
\token{ackermann(1, 2)} by hand, writing each call that the function makes. How
many calls are there in all?

\subsection*{Exercise 9}

Write a program that asks for a date and prints the day of the week it fell on,
such as \textit{That day is a Sunday.} for July 20th, 1969. Reuse the functions of
Listing~\ref{lst:functions:calendar} that you need, and make sure that the day
typed exists in the month typed.

\vfill%
\pagebreak

\forceevenpage{}

\section*{Solutions}
\markright{\MakeUppercase{Solutions}}%
\subsection*{Exercise 1}

The smallest of three numbers is the smallest of the third and of the smaller of
the first two. The version with three parameters calls the one with two, twice.

\lstinputlisting[style=coloredverbatim, caption={Solution for exercise 1, \textit{examples/functions/solutions/smallest.yr}}]{examples/functions/solutions/smallest.yr}

\subsection*{Exercise 2}

\begin{lstlisting}[style=bashVerb]
n: 5
twice: 10 points
four times: 20
sum: 15!
\end{lstlisting}

The first call leaves \token{unit} at its default, the empty string. The third
names \token{value}, so the positional \token{"four times"} goes to the first
free parameter, \token{label}. In the last, \token{unit} is named first, and the
two positional arguments fill \token{label} and \token{value} in order.

\subsection*{Exercise 3}

The body of \token{average} ends with a semicolon, so it has no value.
\token{value} is a parameter, and cannot be assigned: the three cases are
written as one \token{if} used as a value. And \token{low} and \token{high} are
marked \token{@}, so the call must name them. The program prints 5: the average
is 7, above the upper bound.

\begin{lstlisting}[style=coloredverbatim, escapechar=|, caption=Solution for exercise 3]
use std::io;

fn average(a: i32, b: i32)-> i32 {
  |\hcb{(a + b) / 2}|
}

fn clamp(value: i32, @low: i32, @high: i32)-> i32 {
  if value < low {
    |\hcb{low}|
  } else if value > high {
    high
  } else {
    value
  }
}

fn main() {
  let m = average(4, 10);
  println(clamp(m, |\hcb{low-> 0, high-> 5}|));
}
\end{lstlisting}

\subsection*{Exercise 4}

With the base case \token{exponent == 0}, a negative exponent moves away from 0
at each call, and the recursion never ends: the program crashes with a stack
overflow, as in Section~\ref{sec:functions:infinite_recursion}. Testing
\token{exponent <= 0} catches it.

\lstinputlisting[style=coloredverbatim, caption={Solution for exercise 4, \textit{examples/functions/solutions/power.yr}}]{examples/functions/solutions/power.yr}
\vspace{-10pt}%
\begin{lstlisting}[style=bashVerb]
1024
81
1
1
\end{lstlisting}

\subsection*{Exercise 5}

The sum of the digits of \token{n} is its last digit, \token{n \% 10}, plus the
sum of the digits of the others, \token{n / 10}.

\lstinputlisting[style=coloredverbatim, caption={Solution for exercise 5, \textit{examples/functions/solutions/digits.yr}}]{examples/functions/solutions/digits.yr}

\subsection*{Exercise 6}

\lstinputlisting[style=coloredverbatim, caption={Solution for exercise 6, recursive, \textit{examples/functions/solutions/fibonacci.yr}}]{examples/functions/solutions/fibonacci.yr}

\token{fibonacci(30)} answers at once, \token{fibonacci(40)} takes most of a
second, and \token{fibonacci(45)} several seconds. Each call makes two more, and
they compute the same values again and again: \token{fibonacci(40)} calls
\token{fibonacci(39)} and \token{fibonacci(38)}, but \token{fibonacci(39)}
computes \token{fibonacci(38)} a second time, and each of them recomputes all the
smaller ones. In all, \token{fibonacci(40)} makes more than 331 million calls,
and each step of 5 in \token{n} multiplies that by about 11.

The loop keeps the last two numbers of the sequence and moves them forward
\token{n} times, so it makes \token{n} additions in all.

\lstinputlisting[style=coloredverbatim, caption={Solution for exercise 6, with a loop, \textit{examples/functions/solutions/fibonacci\_loop.yr}}]{examples/functions/solutions/fibonacci_loop.yr}
\vspace{-10pt}%
\begin{lstlisting}[style=bashVerb, escapechar=@+]
@+\textcolor{teal!80}{alice@dev:\mtilde}@+$ ./fibonacci_loop
n = 90
fibonacci(90) = 2880067194370816120
\end{lstlisting}

\subsection*{Exercise 7}

\lstinputlisting[style=coloredverbatim, caption={Solution for exercise 7, \textit{examples/functions/solutions/gcd.yr}}]{examples/functions/solutions/gcd.yr}
\vspace{-10pt}%
\begin{lstlisting}[style=bashVerb, escapechar=@+]
@+\textcolor{teal!80}{alice@dev:\mtilde}@+$ ./gcd
a = 1071
b = 462
gcd(1071, 462) = 21
\end{lstlisting}

The remainder of a division by \token{b} is smaller than \token{b}, so for
positive numbers the second argument gets smaller at each call, and eventually
reaches the base case, 0.

\subsection*{Exercise 8}

Each line is one call. The calls with $m = 1$ wait for the value of their inner
call, then make one more call with $m = 0$.

\begin{lstlisting}[style=bashVerb]
ackermann(1, 2) = ackermann(0, ackermann(1, 1))
ackermann(1, 1) = ackermann(0, ackermann(1, 0))
ackermann(1, 0) = ackermann(0, 1)
ackermann(0, 1) = 2
ackermann(0, 2) = 3          so ackermann(1, 1) = 3
ackermann(0, 3) = 4          so ackermann(1, 2) = 4
\end{lstlisting}

There are 6 calls in all, three with $m = 1$ and three with $m = 0$. The value,
4, is the one in the table of Listing~\ref{lst:functions:ackermann}.

\subsection*{Exercise 9}

\token{weekday} and the functions it calls are copied unchanged. The upper bound
of the day depends on the month and the year, so it is computed by
\token{daysIn}, and the numbers from 0 to 6 are turned into names by a
\token{match}, as \token{monthName} does for months.

\lstinputlisting[style=coloredverbatim, caption={Solution for exercise 9, \textit{examples/functions/solutions/weekday.yr}}]{examples/functions/solutions/weekday.yr}
\vspace{-10pt}%
\begin{lstlisting}[style=bashVerb, escapechar=@+]
@+\textcolor{teal!80}{alice@dev:\mtilde}@+$ ./weekday
Year: 1969
Month: 7
Day: 20
That day is a Sunday.
\end{lstlisting}
